% !TEX root = ../main.tex
\section{Introduction}
\label{sec:intro}

Mixture-of-experts (MoE) layers let a language model grow its parameter count without a matching growth in per-token
compute: a router sends each token to a few of many experts, and only those experts run
\citep{shazeer2017outrageously,fedus2022switch}. Qwen3-30B-A3B routes each token to 8 of 128 experts in each of its 48
layers \citep{yang2025qwen3}, and Mixtral and DeepSeek-V3 follow the same pattern \citep{jiang2024mixtral,deepseek2024v3}.
During prefill a serving system processes the prompt in chunks of $T$ tokens \citep{agrawal2024sarathi}, and the router
distributes each chunk's $8T$ assignments unevenly over the experts, differently in every layer and for every input.
Efficient execution should therefore give different experts different amounts of work.

Server inference accelerators make that hard. Qualcomm Cloud AI 100 \citep{chatha2021qualcomm}, like XLA-compiled TPU
programs \citep{lepikhin2021gshard} and the statically scheduled Groq TSP \citep{abts2020tsp}, runs a model as a program
compiled ahead of time: tensor shapes, the partition of every operator over cores and devices, and the schedule are fixed
before the first request, and the compiler rejects control flow that depends on data. Expert work must be expressed
through compiled shapes. A uniform capacity of $T$ rows per expert accepts any routing, but it executes padded rows,
16$\times$ the routed assignments for this model, and the compiler skips them only in whole row tiles. Capacities fitted
to observed loads do not escape this, because which experts are busy changes with the workload (\S\ref{sec:mot-pad}).

Reducing expert work is only half of the problem. Expert outputs must return to their tokens and be summed over each
token's experts, and the shapes and stages that capacity savings introduce change how that aggregation executes. In a
controlled experiment that holds everything else fixed, ranked capacities cut the full-model latency of 512-token chunks
by 27\% with \sys's aggregation but by only 4\% with a repaired dense combine, and at 128 tokens they make the dense
version 37\% slower (\S\ref{sec:mot-red}). Choosing capacities to minimize padded rows is an incomplete objective:
capacities and aggregation have to be designed together.

Efficient aggregation in turn depends on how the graph expresses it. The mathematics is standard: each device sums the
rows of the experts it computed into per-token partials, and the partials are summed across devices, as GPU systems do
with sparse kernels \citep{hwang2023tutel,vllmmoe}. Through a graph compiler, however, equivalent expressions of that
computation execute very differently. Three graphs that produce bit-identical logits differ only in how they tell the
compiler which device owns which rows: global row indices and explicit per-device slices both make it ship expert outputs
between devices, 110$\times$ the traffic, and run 3.4--5.3$\times$ slower than a layout that exposes the device as a batch
axis of the gather (\S\ref{sec:mot-own}). Within a device, the reduction must also stay on the core that owns each token:
summing a token's rows across cores costs 11--15\%.

We present \sys, a design for MoE prefill that keeps the program static and moves the decisions it cannot compile into
data. Each SoC ranks its 32 experts by load at run time and assigns them to fixed-capacity tiers by rank; the compiled
program fixes the tiers, their lanes and capacities, and the run time decides which expert each lane serves. Capacities
bound to ranks can be small because the order statistics of per-device load are bounded and stable across workloads,
while the identities of the busy experts are not (Figure~\ref{fig:rank}). Expert outputs are aggregated in a layout that
preserves device ownership and assigns each token's sum to one core, and that does not depend on the tiers' shapes. Both
parts need only static shapes, a sort and gathers with run-time indices.

We implement \sys as graph transformations of the exported Qwen3-30B-A3B and evaluate all 48 layers with MXFP6 expert
weights on the four SoCs of one Cloud AI 100 Ultra card, with the KV cache on device. \sys prefills 128-, 256- and
512-token chunks in 105.3, 141.5 and 225.3\,ms, 4.74$\times$, 3.65$\times$ and 4.68$\times$ faster than the production compilation and
2.37$\times$, 3.55$\times$ and 4.56$\times$ faster than an expert-parallel program with the exported combine. With
ranking and weight indirection matched, the ranked capacities give 1.10--1.37$\times$; fixed-size blocks under the same
compiler are 1.26--1.71$\times$ slower. Capacities calibrated in advance overflowed in none of 1872 evaluation chunks of
eight workloads, and a forced overflow is recovered bit-exactly. We separate what transfers to other statically
partitioned targets from the behavior of this compiler (\S\ref{sec:transfer}).

The paper makes three contributions:
\begin{itemize}
\item A characterization of MoE prefill under static partitioning: padded expert rows are executed work, capacity savings
  pay only with efficient aggregation, and graphs with identical arithmetic place that aggregation very differently
  (\S\ref{sec:mot}).
\item A design that adapts expert work and preserves the aggregation structure: capacities bound to per-device load ranks,
  with exact overflow recovery, and token aggregation that keeps device ownership and token-local reduction through
  compilation (\S\ref{sec:design}, \S\ref{sec:impl}).
\item A full-model evaluation that establishes the capacity--aggregation interaction with controlled ablations,
  equivalent-arithmetic placement experiments and an attributed optimization ladder (\S\ref{sec:eval}).
\end{itemize}
